\documentclass[12pt,a4paper]{article}

\usepackage[utf8]{inputenc}
\usepackage[T1]{fontenc}
\usepackage[brazilian]{babel}
\usepackage{amsmath}
\usepackage{amssymb}
\usepackage{mathtools}
\usepackage[margin=2.5cm]{geometry}
\usepackage{booktabs}
\usepackage{array}
\usepackage{enumitem}
\usepackage[colorlinks=true,linkcolor=blue,citecolor=blue,urlcolor=blue]{hyperref}

\title{Fundamentação Teórica e Fórmulas de Escoragem\\[2pt]
\large Índices Afya de Qualidade de Vida: MedQoL (médicos) e Afya MedQoL Student (estudantes de medicina)}
\author{Afya}
\date{\today}

\begin{document}

\maketitle

\begin{abstract}
\noindent
Este documento consolida, em formato de metodologia de artigo científico, as
fórmulas teóricas utilizadas na biblioteca \texttt{afya\_medqol} para o
escoramento de dois instrumentos de qualidade de vida: o \textbf{Afya MedQoL Physician}
(médicos formados, 13 itens, calibração 2024\_2) e o \textbf{Afya MedQoL Student} (estudantes
de medicina, 8 itens). Ambos os instrumentos utilizam modelos de Teoria de
Resposta ao Item (TRI) da família do Modelo de Resposta Gradual de Samejima
(\emph{Graded Response Model}, GRM), com escoragem bayesiana via Esperança a
Posteriori (\emph{Expected A Posteriori}, EAP). Ambos os instrumentos
reportam o escore em $\theta$ e em T-score. As subseções
\ref{sec:medqol} e \ref{sec:iqol}, dentro da Metodologia, são
independentes e podem ser lidas separadamente.
\end{abstract}

\tableofcontents

% =====================================================================
\section{Introdução}
% =====================================================================

Os dois instrumentos aqui descritos diferem na estrutura de fatores latentes
que representam a qualidade de vida do respondente:

\begin{itemize}
    \item O \textbf{Afya MedQoL Physician} assume \textbf{três domínios ortogonais e
    independentes} ($F_1$, $F_2$, $F_3$), cada um escorado por um subconjunto
    de itens que carrega em um único fator (estrutura simples).
    \item O \textbf{Afya MedQoL Student} assume uma estrutura \textbf{bifatorial}: cada item
    carrega simultaneamente em um fator geral de qualidade de vida
    ($\theta_G$, comum aos 8 itens) e em um fator específico do seu domínio
    ($\theta_{S}$), sendo os fatores especificos ortogonais entre si e ao
    fator geral.
\end{itemize}

Em ambos os casos, a escoragem é \textbf{100\% determinística}: a grade de
quadratura numérica é fixa (sem amostragem estocástica e sem semente
aleatória), de modo que os mesmos parâmetros calibrados e o mesmo banco de
respostas produzem sempre os mesmos escores, em qualquer máquina.

% =====================================================================
\section{Metodologia}
\label{sec:metodologia}
% =====================================================================

Esta seção apresenta, em duas subseções independentes, o modelo de resposta
ao item, o procedimento de escoragem bayesiana e as conversões de escala
utilizados para cada instrumento.

\subsection{Afya MedQoL Physician --- Médicos formados}
\label{sec:medqol}

\subsubsection{Estrutura do instrumento}

O instrumento é composto por 13 itens distribuídos em três domínios
independentes:

\begin{table}[h]
\centering
\begin{tabular}{@{}llc@{}}
\toprule
Domínio & Construto & N.º de itens \\
\midrule
$F_1$ & Qualidade de Vida & 6 \\
$F_2$ & Suporte Institucional / Percepção do Trabalho & 4 \\
$F_3$ & Estresse Percebido & 3 \\
\bottomrule
\end{tabular}
\end{table}

Cada item $j$ é respondido em uma escala ordinal de $K=5$ categorias
($Y_j \in \{1,\dots,5\}$) e carrega em exatamente um domínio $f(j) \in
\{F_1,F_2,F_3\}$ (estrutura simples).

\subsubsection{Modelo de resposta ao item: GRM de Samejima}

Seja $\boldsymbol{\theta} = (\theta_{F_1}, \theta_{F_2}, \theta_{F_3})^\top$ o
vetor de traços latentes do respondente, com distribuição \emph{a priori}

\begin{equation}
\boldsymbol{\theta} \sim \mathcal{N}_3(\mathbf{0}, \Sigma), \qquad
\Sigma = I_3
\label{eq:medqol-prior}
\end{equation}

isto é, os três domínios são mutuamente independentes por construção
($\Sigma$ é a matriz identidade). Para o item $j$, com discriminação $a_j$ e
limiares (\emph{thresholds}) $b_{j,1} < b_{j,2} < b_{j,3} < b_{j,4}$, o Modelo
de Resposta Gradual de Samejima~\cite{samejima1969} define a probabilidade
acumulada de resposta em categoria igual ou superior a $k$:

\begin{equation}
P^*_{jk}(\boldsymbol{\theta}) \;=\; P(Y_j \geq k \mid \boldsymbol{\theta})
\;=\; \frac{1}{1 + \exp\!\big(-a_j\,(\theta_{f(j)} - b_{j,k-1})\big)},
\qquad k = 2, \dots, 5
\label{eq:medqol-cumulative}
\end{equation}

com as condições de contorno $P^*_{j1}(\boldsymbol{\theta}) \equiv 1$ e
$P^*_{j,6}(\boldsymbol{\theta}) \equiv 0$. A probabilidade de resposta exatamente
na categoria $k$ é obtida por diferença:

\begin{equation}
P_{jk}(\boldsymbol{\theta}) \;=\; P^*_{jk}(\boldsymbol{\theta}) - P^*_{j,k+1}(\boldsymbol{\theta}),
\qquad k = 1, \dots, 5
\label{eq:medqol-category}
\end{equation}

Como cada item carrega em um único domínio, $\theta_{f(j)}$ denota a
coordenada de $\boldsymbol{\theta}$ correspondente ao domínio do item $j$.

\subsubsection{Escoragem bayesiana: Esperança a Posteriori (EAP)}

Dado o vetor de respostas $\mathbf{Y} = (Y_1, \dots, Y_{13})$ de um
respondente, a log-verossimilhança é

\begin{equation}
\ell(\boldsymbol{\theta} \mid \mathbf{Y}) \;=\; \sum_{j=1}^{13} \log P_{j,Y_j}(\boldsymbol{\theta})
\label{eq:medqol-loglik}
\end{equation}

somando apenas sobre os itens efetivamente respondidos (itens omissos ou com
código de não resposta são excluídos do somatório). A densidade \emph{a
posteriori} é proporcional ao produto da verossimilhança pelo prior
trivariado~\eqref{eq:medqol-prior}:

\begin{equation}
p(\boldsymbol{\theta} \mid \mathbf{Y}) \;\propto\;
\exp\!\big(\ell(\boldsymbol{\theta}\mid\mathbf{Y})\big) \cdot
\varphi_3(\boldsymbol{\theta}; \mathbf{0}, \Sigma)
\label{eq:medqol-posterior}
\end{equation}

onde $\varphi_3(\cdot\,;\mathbf{0},\Sigma)$ é a densidade normal trivariada.
Como $\Sigma = I_3$ é diagonal e a estrutura de itens é simples, a
integração conjunta em $\mathbb{R}^3$ é matematicamente equivalente a três
integrações unidimensionais independentes — a "régua independente" por
domínio. Na implementação, a integração é feita numericamente por
quadratura fixa sobre uma grade cúbica
$\{\boldsymbol{\theta}^{(1)}, \dots, \boldsymbol{\theta}^{(G)}\}$ em
$[-5,5]^3$ ($G = 25^3$ nós), com pesos $w^{(g)} \propto
\varphi_3(\boldsymbol{\theta}^{(g)};\mathbf{0},\Sigma)$ normalizados para
somar 1. O escore EAP é a média \emph{a posteriori}:

\begin{equation}
\hat{\boldsymbol{\theta}} \;=\; \mathbb{E}[\boldsymbol{\theta} \mid \mathbf{Y}]
\;\approx\; \sum_{g=1}^{G} \boldsymbol{\theta}^{(g)} \cdot
\frac{w^{(g)} \exp\big(\ell(\boldsymbol{\theta}^{(g)}\mid\mathbf{Y})\big)}
{\sum_{g'=1}^{G} w^{(g')} \exp\big(\ell(\boldsymbol{\theta}^{(g')}\mid\mathbf{Y})\big)}
\label{eq:medqol-eap}
\end{equation}

\subsubsection{Escore composto global}

O escore global combina os três domínios, com o domínio $F_3$ (Estresse
Percebido) invertido apenas nesta composição — quanto maior o estresse, menor
a contribuição para o bem-estar global:

\begin{equation}
\hat\theta_{\text{global}} \;=\; w_{F_1}\,\hat\theta_{F_1} \;+\; w_{F_2}\,\hat\theta_{F_2} \;-\; w_{F_3}\,\hat\theta_{F_3}
\label{eq:medqol-global-theta}
\end{equation}

com pesos proporcionais à discriminação total de cada domínio na
calibração:

\begin{equation}
w_f \;=\; \frac{\sum_{j : f(j)=f} |a_j|}{\sum_{j=1}^{13} |a_j|}, \qquad f \in \{F_1,F_2,F_3\}
\label{eq:medqol-weights}
\end{equation}

\subsubsection{T-score}

Por compatibilidade com a literatura de TRI, também é reportado um T-score
por domínio e global, com $\theta \sim \mathcal{N}(0,1)$ \emph{a priori}:

\begin{equation}
T_f = 50 + 10\,\hat\theta_f, \qquad
T_{\text{global}} = 50 + 10\,\hat\theta_{\text{global}}
\label{eq:medqol-tscore}
\end{equation}

\subsection{Afya MedQoL Student --- Estudantes de medicina}
\label{sec:iqol}

\subsubsection{Estrutura do instrumento}

O Afya MedQoL Student (\emph{Student Quality of Life Index}) é composto por 8 itens
distribuídos em três domínios específicos, todos carregando também em um
fator geral comum:

\begin{table}[h]
\centering
\begin{tabular}{@{}cllc@{}}
\toprule
Domínio & Construto & Itens & Peso global ($w_f$) \\
\midrule
1 & Bem-estar psicológico & \texttt{F1\_1\_overallqol}, \texttt{F1\_2\_satisfactionwithhealth}, \texttt{F1\_3\_enjoymentoflife}, \texttt{F1\_4\_perceivedmeaninginlife} & 0,494 \\
2 & Vitalidade & \texttt{F2\_1\_energyfordailyactivities}, \texttt{F2\_2\_satisfactionwithsleep} & 0,172 \\
3 & Capacidade funcional percebida & \texttt{F3\_1\_performdailyactivities}, \texttt{F3\_2\_capacityforwork} & 0,335 \\
\bottomrule
\end{tabular}
\end{table}

Cada item $j$ é respondido em escala ordinal de $K=5$ categorias e nenhum
item é invertido (todos unidirecionais).

\subsubsection{Modelo de resposta ao item: GRM bifatorial}

Diferentemente do MedQoL, o Afya MedQoL Student modela cada respondente por um vetor
$(\theta_G, \theta_{S_1}, \theta_{S_2}, \theta_{S_3})$, em que $\theta_G$ é um
fator \textbf{geral} de qualidade de vida comum aos 8 itens, e $\theta_{S_f}$
é o fator \textbf{específico} do domínio $f$. Todos os fatores são
mutuamente ortogonais, com prior padronizado independente:

\begin{equation}
\theta_G \sim \mathcal{N}(0,1), \qquad
\theta_{S_f} \sim \mathcal{N}(0,1) \; \perp \; \theta_G, \quad f=1,2,3
\label{eq:iqol-prior}
\end{equation}

Para o item $j$, com discriminação geral $a_{G,j}$, discriminação específica
$a_{S,j}$ (no domínio $f(j)$) e interceptos (parametrização \emph{mirt})
$d_{j,1} > d_{j,2} > d_{j,3} > d_{j,4}$, o modelo bifatorial de Samejima
(intercept parameterization) define:

\begin{equation}
P^*_{jk}(\theta_G, \theta_{S_{f(j)}}) \;=\; P(Y_j \geq k \mid \theta_G,\theta_{S_{f(j)}})
\;=\; \frac{1}{1 + \exp\!\big(-(a_{G,j}\,\theta_G + a_{S,j}\,\theta_{S_{f(j)}} + d_{j,k-1})\big)},
\quad k=2,\dots,5
\label{eq:iqol-cumulative}
\end{equation}

com $P^*_{j1} \equiv 1$ e $P^*_{j,6} \equiv 0$, e probabilidade de categoria

\begin{equation}
P_{jk}(\theta_G, \theta_{S_{f(j)}}) \;=\; P^*_{jk}(\theta_G, \theta_{S_{f(j)}}) - P^*_{j,k+1}(\theta_G, \theta_{S_{f(j)}})
\label{eq:iqol-category}
\end{equation}

\subsubsection{Escoragem bayesiana bifatorial (EAP por integração numérica)}

A escoragem bifatorial segue o procedimento clássico de
Gibbons \& Hedeker~\cite{gibbonshedeker1992}: o fator geral $\theta_G$ é
integrado \textbf{conjuntamente} sobre os três domínios, e cada fator
específico $\theta_{S_f}$ é obtido por integração condicional a $\theta_G$.

Para cada domínio $f$, seja $\mathbf{Y}_f$ o subvetor de respostas aos itens
de $f$. Define-se, para cada valor de $\theta_G$ na grade de quadratura, as
integrais marginais sobre $\theta_{S_f}$ (com $\phi(\cdot)$ a densidade
normal padrão):

\begin{equation}
L_f(\theta_G) \;=\; \int_{-\infty}^{\infty}
\left[\prod_{j \,:\, f(j)=f} P_{j,Y_j}(\theta_G,\theta_{S_f})\right]
\phi(\theta_{S_f}) \; d\theta_{S_f}
\label{eq:iqol-Lf}
\end{equation}

\begin{equation}
M_f(\theta_G) \;=\; \int_{-\infty}^{\infty} \theta_{S_f}
\left[\prod_{j \,:\, f(j)=f} P_{j,Y_j}(\theta_G,\theta_{S_f})\right]
\phi(\theta_{S_f}) \; d\theta_{S_f}
\label{eq:iqol-Mf}
\end{equation}

Ambas as integrais são aproximadas por quadratura fixa sobre uma grade de
121 nós equiespaçados em $[-6,6]$, com pesos $\propto \phi(\theta_{S_f}^{(g)})$
normalizados. A posterior marginal do fator geral, integrando os três
domínios conjuntamente, é

\begin{equation}
p(\theta_G \mid \mathbf{Y}) \;\propto\; \phi(\theta_G) \prod_{f=1}^{3} L_f(\theta_G)
\label{eq:iqol-postG}
\end{equation}

também avaliada sobre a mesma grade de 121 nós e normalizada para integrar 1.
O EAP de cada fator específico é obtido integrando a esperança condicional
$\mathbb{E}[\theta_{S_f} \mid \theta_G] = M_f(\theta_G)/L_f(\theta_G)$ sobre
a posterior marginal de $\theta_G$:

\begin{equation}
\hat\theta_{S_f} \;=\; \int_{-\infty}^{\infty}
\frac{M_f(\theta_G)}{L_f(\theta_G)} \; p(\theta_G \mid \mathbf{Y}) \; d\theta_G,
\qquad f = 1,2,3
\label{eq:iqol-eap}
\end{equation}

novamente aproximada pela mesma grade fixa de 121 nós em $[-6,6]$, garantindo
determinismo total do procedimento.

\subsubsection{Escore composto global}

O escore global do Afya MedQoL Student é o composto ponderado dos três fatores específicos,
\textbf{sem} inversão de sinal (nenhum domínio é invertido no Afya MedQoL Student):

\begin{equation}
\hat\theta_{\text{global}} \;=\; \sum_{f=1}^{3} w_f \, \hat\theta_{S_f}
\label{eq:iqol-global-theta}
\end{equation}

com pesos $w_f$ fixos, proporcionais à discriminação relativa de cada
domínio na calibração (Tabela na Subseção~\ref{sec:iqol}: $w_1=0{,}494$,
$w_2=0{,}172$, $w_3=0{,}335$).

\subsubsection{T-score}

Diferentemente do MedQoL, $\hat\theta_{\text{global}}$ do Afya MedQoL Student é um composto
ponderado e, portanto, não tem variância unitária \emph{a priori}. O T-score
é obtido por padronização empírica, usando a média $\mu_g$ e o desvio-padrão
$\sigma_g$ de $\hat\theta_{\text{global}}$ na amostra de referência da
calibração:

\begin{equation}
z = \frac{\hat\theta_{\text{global}} - \mu_g}{\sigma_g}, \qquad
T_{\text{global}} = 50 + 10\,z
\label{eq:iqol-tscore}
\end{equation}

% =====================================================================
\section{Notas de reprodutibilidade}
% =====================================================================

Em ambos os instrumentos, os parâmetros de calibração (discriminações,
limiares/interceptos, pesos e, no caso do Afya MedQoL Student, também $\mu_g/\sigma_g$) são
\textbf{fixos e congelados} a partir dos respectivos artigos de validação
psicométrica, e não são reestimados a cada aplicação.
A combinação (parâmetros fixos + quadratura fixa + código determinístico)
garante que qualquer banco de respostas, processado em qualquer máquina,
produza exatamente os mesmos escores.

% =====================================================================
\begin{thebibliography}{9}

\bibitem{samejima1969}
Samejima F.
\emph{Estimation of latent ability using a response pattern of graded scores}.
Psychometrika Monograph Supplement, 1969.

\bibitem{stockinglord1983}
Stocking ML, Lord FM.
\emph{Developing a common metric in item response theory}.
Applied Psychological Measurement. 1983;7(2):201--210.

\bibitem{gibbonshedeker1992}
Gibbons RD, Hedeker DR.
\emph{Full-information item bi-factor analysis}.
Psychometrika. 1992;57(3):423--436.

\bibitem{gobbo2025medqol}
Gobbo Jr M, et al.
\emph{[Desenvolvimento e validação psicométrica do índice Afya MedQoL]}.
BMJ Open. 2025;15:e102783.

\bibitem{gobbo2026iqol}
Gobbo M Jr, Silveira de Resende M, Santos IM, Cardoso de Moura E, Almeida Pedro R.
\emph{Development and psychometric validation of the 8-item Student Quality
of Life Index (IQoL) using item response theory (IRT)}.
BMJ Open. 2026;16:e106371.

\bibitem{kolenbrennan2014}
Kolen MJ, Brennan RL.
\emph{Test Equating, Scaling, and Linking: Methods and Practices}. 3rd ed.
Springer, 2014.

\bibitem{cella2010}
Cella D, et al.
\emph{The Patient-Reported Outcomes Measurement Information System (PROMIS):
progress of an NIH Roadmap cooperative group during its first two years}.
Medical Care. 2007;45(5 Suppl 1):S3--S11.

\end{thebibliography}

\end{document}
